\section{Reads at two instants}
\label{sec:time}

\subsection{What each coordinate asks}

\shipped{} Every operation that resolves takes \texttt{as\_of}, the valid
instant $\tv$, and \texttt{as\_known}, the transaction instant $\tk$; either
absent is the server clock. The raw read, \op{GET /v1/graph}, takes the same two
bounds and resolves nothing: with neither, it returns every row matching its
filter, and it reports when its row limit stopped it.

Holding $\tk$ at the present and moving $\tv$ asks about the world as currently
known: who held a position in 1998. Holding $\tv$ and moving $\tk$ asks what the
store would have answered at an earlier moment, which is what makes a past
answer reproducible and what an audit of a decision needs. Moving both asks what
was believed on one day about another day. Table~\ref{tab:example} shows the
three on one pair.

\begin{table}[h]
\centering
\small
\begin{tabular}{@{}clllll@{}}
\toprule
& value & \texttt{at} & \texttt{until} & filed (knowable) & kind \\
\midrule
$a_1$ & \texttt{person/x} & 1993-01-01 & 1998-01-01 & 2026-01-10 & edge \\
$a_2$ & \texttt{person/y} & 1997-01-01 & 2002-01-01 & 2026-01-10 & edge \\
$a_3$ & \texttt{person/x} & 1993-01-01 & 1997-01-01 & 2026-03-02 & edge, a version of $a_1$ \\
\midrule
\multicolumn{6}{@{}l}{\emph{Resolved on the pair (\texttt{position/p}, \texttt{held\_by}), undeclared and so holding several values:}} \\
\multicolumn{6}{@{}l}{\quad $(\tv, \tk) = (\text{1997-07-01}, \text{2026-02-01})$: holds $a_2$ and $a_1$} \\
\multicolumn{6}{@{}l}{\quad $(\tv, \tk) = (\text{1997-07-01}, \text{2026-04-01})$: holds $a_2$; $a_3$ replaced $a_1$ and ended in 1997} \\
\multicolumn{6}{@{}l}{\quad $(\tv, \tk) = (\text{1995-07-01}, \text{2026-04-01})$: holds $a_3$} \\
\bottomrule
\end{tabular}
\caption{An illustration, not data. $a_3$ corrects the end of $a_1$'s term. As
known before the correction, the old account holds; as known after it, the new
one does; asking about another year changes neither. The first line is the case
the previous version could not answer: both terms hold in 1997, and neither ends
the other.}
\label{tab:example}
\end{table}

\subsection{The previous rule, in these terms}

The version measured in the first CronQuestions run read at one instant,
\texttt{as\_of}, and bounded the knowable instant by it. In the terms of this
section that is a read at $(\tv, \tk) = (\infty, T)$: transaction time only,
with every statement ever filed eligible regardless of when it was so. A
question about 1998 asked at $T = $ 1998 found the rows the store had heard of
by 1998, and history filed in 2026 was not among them. The measurement recorded
this as zero correct answers on point-in-time questions when the whole
knowledge graph was filed at one instant (\S\ref{sec:eval-first}). A system that
reads by transaction time alone has this behaviour by construction; Datomic's
\texttt{as-of} database value, for instance, is a transaction-time
filter~\cite{datomic}, and valid time is left to be modelled in the schema.

The guard that produced the behaviour, Proposition~\ref{prop:noback}, was kept.
What changed is that the question about the world now has its own coordinate,
so the guard applies to the coordinate it was meant for.

\subsection{When filing time does not matter}

The CronQuestions lane files the whole knowledge graph twice, once with every
row written at one server instant and once with each row's write instant set to
its own $\mathit{at}$, and asks every question at $\tk$ after the load. The
following property says when the two must agree.

\begin{proposition}[Valid-time answers independent of filing time]
\label{prop:filing}
Take a pair that holds several values and contains no retraction, and suppose
that all versions of each of its statements share one knowable instant. Then for
every $\tk$ at or after the latest knowable instant among its rows, the pair
holds at $(\tv, \tk)$ exactly its rows that are open at $\tv$, whatever the
knowable instants are.
\end{proposition}
\begin{proof}
Every row is in $K_{\tk}$. Because the versions of each statement share their
knowable instant, every version is current. The several-value rule keeps each
non-retraction that is open at $\tv$ and not ended by a retraction, and the pair
has none. No step compares knowable instants.
\end{proof}

For a pair holding one value, the knowable instant enters through $\succ$ only
when two statements share an $\mathit{at}$; if none do, the same conclusion
holds with ``the latest statement begun by $\tv$, if open'' in place of ``its
rows open at $\tv$''. Every relation in the lane is an edge and undeclared, so
every pair holds several values, and the lane files no retraction. The knowledge
graph does repeat statements: \res{cron/dupkeys} (subject, relation, object,
start year) keys occur on more than one line, \res{cron/dupends} of them with
different end years. Under the first clock every row shares one write instant,
and under the second the versions of a statement share $\mathit{at}$, so they
share a knowable instant under both. The proposition therefore predicts that
the two loads answer identically, and they do, question for question
(\S\ref{sec:eval-now}).

\subsection{Diff between two points}

\shipped{} \op{diff} takes two points $a = (\tv^a, \tk^a)$ and
$b = (\tv^b, \tk^b)$ with $b$ no earlier than $a$ on either coordinate and
$a \ne b$. Only a pair with a statement that began or ended in
$(\tv^a, \tv^b]$, or became knowable in $(\tk^a, \tk^b]$, can hold something
different at the two points, so the operation selects those pairs with one
query, resolves each at both points, and reports three kinds of change:
a value held at $b$ and not at $a$ (\emph{asserted}); for a one-valued pair, a
value replaced by another (\emph{superseded}, reported once however many
intermediate values there were); and a value held at $a$ and not at $b$
(\emph{retracted}), together with what ended it --- the retraction, the later
value, or the statement's own $\until$.

Because the report is computed from what each pair holds at the two points and
not from the rows in between, a version that loses the order produces no change,
and one that repeats the value already held is reported as nothing. Holding the
transaction coordinates equal gives what happened in the world between two
valid instants as currently known; holding the valid coordinates equal gives
what the store learned between two transaction instants, which is an audit of
the record. The number of pairs examined is bounded, 10{,}000 by default, and
the response says when the bound or a pair's row limit was reached.
